\BeleskeID{09_2-s033}
% element: 09_2-s033:e04
Polazna ravnoteža struja uključuje efekat kratkog kanala:
\[*1.\quad k_p\frac1{1+\frac{V_I-V_{DD}-V_{Tp}}{L_pE_{Cp}}}(V_I-V_{DD}-V_{Tp})^2=k_n\frac1{1+\frac{V_O}{L_nE_{Cn}}}\left(2V_O(V_I-V_{Tn})-V_O^2\right)\]
Njeno neposredno diferenciranje dalo bi složen izraz. Očekuje se mali izlazni napon $V_{O(IH)}$, blizu nule, i ulazni napon $V_{IH}$ blizu $V_{DD}+V_{Tp}$. Zato se u izvođenju uvode prikazane pretpostavke o malim količnicima i imenitelji aproksimiraju jedinicom. Pojednostavljena ravnoteža struja diferencira se uz $\partial V_O/\partial V_I=-1$, a dobijena linearna jednačina sređuje po ulaznom naponu. Zatim se rešava zajedno sa ravnotežom struja.
